% ===================== Fig. 3 : multi-graph fusion =====================
\begin{figure*}[t]
\centering
\resizebox{\textwidth}{!}{%
\definecolor{boxcol}{RGB}{52,86,122}%
\begin{tikzpicture}[
  font=\small, x=1cm, y=1cm,
  >={Stealth[length=3.6mm, width=2.6mm]},
  rg/.style={circle, draw=black!65, fill=white, inner sep=0pt, minimum size=7mm, font=\small},
  foc/.style={rg, draw=black, line width=1.1pt, fill=yellow!25},
  blk/.style={draw=black!70, rounded corners=2pt, align=center, inner sep=4pt, fill=white},
  op/.style={blk, fill=black!5, font=\small},
  cell/.style={draw=black!45, minimum size=6.4mm, inner sep=0pt, font=\normalsize},
  f/.style={->, draw=black!65, line width=1.2pt},
  F/.style={->, draw=black!85, line width=1.9pt},
  lbl/.style={font=\normalsize, text=black!75, inner sep=2pt},
  blbl/.style={font=\large, text=black!80, inner sep=2pt},
  wt/.style={font=\large, inner sep=2pt},
  grp/.style={draw=boxcol, dashed, dash pattern=on 2.2pt off 1.6pt, rounded corners=4pt,
              inner sep=7pt, line width=1.3pt}
]

% five regions, identical layout in every panel; i is the focal region
\def\nodelayout{
  \node[foc] (n1) at (0.00,1.35) {$i$};
  \node[rg]  (n2) at (1.35,1.68) {$j$};
  \node[rg]  (n3) at (1.82,0.45) {$k$};
  \node[rg]  (n4) at (0.60,0.03) {$l$};
  \node[rg]  (n5) at (-0.60,0.60) {$m$};
}
% row of A_{i,:} as five shaded cells: #1 x-origin, #2..#6 fills
\newcommand{\rowstrip}[7]{%
  \node[blbl, anchor=south] at (#1+1.1,-2.32) {#7};
  \foreach \c [count=\q from 0] in {#2,#3,#4,#5,#6}
    \node[cell, fill=\c] at (#1+0.55*\q,-2.75) {};
  \foreach \t [count=\q from 0] in {i,j,k,l,m}
    \node[font=\normalsize, text=black!60] at (#1+0.55*\q,-3.25) {$\t$};
}

%% ---------------- (a) the four adjacencies ----------------
\begin{scope}[shift={(0,0)}]
  \nodelayout
  \draw[teal!70, line width=1.0pt] (n1)--(n2) (n1)--(n5) (n2)--(n3) (n3)--(n4) (n4)--(n5) (n1)--(n4);
  \node[lbl, anchor=north, text width=46mm, align=center] (la) at (0.6,-0.45) {$A^{\mathrm{geo}}$: distance kernel\\[-1pt]\textcolor{teal!70}{fixed, 6-NN sparsified}};
  \begin{scope}[on background layer]\node[grp, fit=(n1)(n2)(n3)(n4)(n5)(la)] {};\end{scope}
\end{scope}
\begin{scope}[shift={(6.1,0)}]
  \nodelayout
  \draw[blue!65, line width=1.1pt, ->, shorten >=2.5pt, shorten <=2.5pt] (n5) to (n1);
  \draw[blue!65, line width=1.4pt, ->, shorten >=2.5pt, shorten <=2.5pt] (n1) to (n2);
  \draw[blue!65, line width=1.1pt, ->, shorten >=2.5pt, shorten <=2.5pt] (n4) to (n3);
  \draw[blue!65, line width=1.1pt, ->, shorten >=2.5pt, shorten <=2.5pt] (n3) to (n2);
  \node[lbl, anchor=north, text width=48mm, align=center] (lb) at (0.6,-0.45) {$A^{\mathrm{adp}}$: learned, static\\[-1pt]\textcolor{blue!65}{$\mathrm{softmax}(\mathrm{ReLU}(E_1E_2^{\!\top}))$}};
  \begin{scope}[on background layer]\node[grp, fit=(n1)(n2)(n3)(n4)(n5)(lb)] {};\end{scope}
\end{scope}
\begin{scope}[shift={(12.2,0)}]
  \nodelayout
  \draw[orange!80, line width=0.9pt, ->, shorten >=2.5pt, shorten <=2.5pt] (n2) to (n1);
  \draw[orange!80, line width=1.6pt, ->, shorten >=2.5pt, shorten <=2.5pt] (n4) to (n1);
  \draw[orange!80, line width=1.2pt, ->, shorten >=2.5pt, shorten <=2.5pt] (n5) to (n1);
  \draw[orange!80, line width=0.9pt, ->, shorten >=2.5pt, shorten <=2.5pt] (n3) to (n2);
  \node[lbl, anchor=north, text width=50mm, align=center] (lc) at (0.6,-0.45) {$A^{\mathrm{dyn}}$: attention on $\mathbf{h}$\\[-1pt]\textcolor{orange!85}{changes every forecast origin}};
  \begin{scope}[on background layer]\node[grp, fit=(n1)(n2)(n3)(n4)(n5)(lc)] {};\end{scope}
\end{scope}
\begin{scope}[shift={(18.3,0)}]
  \nodelayout
  \draw[purple!70, line width=2.2pt] (n1)--(n3);
  \draw[purple!70, line width=1.2pt] (n2)--(n4);
  \draw[purple!70, line width=0.8pt] (n1)--(n2);
  \node[rg, fill=black!10, draw=black!35] at (-0.60,0.60) {$m$};
  \node[lbl, anchor=north, text width=50mm, align=center] (ld) at (0.6,-0.45) {$A^{\mathrm{gen}}$: genomic linkage\\[-1pt]\textcolor{purple!70}{genomes before the cut-off}};
  \begin{scope}[on background layer]\node[grp, fit=(n1)(n2)(n3)(n4)(n5)(ld)] {};\end{scope}
\end{scope}

%% ---------------- (b) row i of each adjacency, weighted by the gate ----------------
\rowstrip{-0.55}{white}{teal!55}{white}{teal!35}{teal!55}{$A^{\mathrm{geo}}_{i,:}$}
\rowstrip{5.55}{white}{white}{blue!50}{white}{white}{$A^{\mathrm{adp}}_{i,:}$}
\rowstrip{11.65}{white}{orange!40}{white}{orange!70}{white}{$A^{\mathrm{dyn}}_{i,:}$}
\rowstrip{17.75}{white}{purple!30}{purple!75}{white}{black!12}{$A^{\mathrm{gen}}_{i,:}$}


\node[wt, text=teal!60!black]   at (1.55,-3.95) {$\times\,\pi_{i,\mathrm{geo}}$};
\node[wt, text=blue!65]         at (7.65,-3.95) {$\times\,\pi_{i,\mathrm{adp}}$};
\node[wt, text=orange!85!black] at (13.75,-3.95) {$\times\,\pi_{i,\mathrm{dyn}}$};
\node[wt, text=purple!70]       at (19.85,-3.95) {$\times\,\pi_{i,\mathrm{gen}}$};

\foreach \x in {1.55,7.65,13.75,19.85} \draw[f] (\x,-3.4) -- (\x,-3.75);
\foreach \x in {1.55,7.65,13.75,19.85} \draw[f, -] (\x,-4.15) -- (\x,-4.85);


%% ---------------- (c) the fused row, gate and message passing ----------------
% all three blocks share one width and one centre line at x = 26.0
\node[op, anchor=north, text width=54mm] (gate) at (26.0,2.00)
  {per-region gate $\pi_i=\mathrm{softmax}(W_\pi\mathbf{h}_i)\in\Delta^{3}$};
\node[op, anchor=north, text width=54mm] (kl) at (26.0,0.70)
  {KL: $A^{\mathrm{gen}}_{i,:}\,\|\,A^{\mathrm{adp}}_{i,:}$ \ (Eq.~8)};
\node[op, anchor=north, text width=54mm] (mp) at (26.0,-0.95)
  {$\mathbf{g}_i\!\leftarrow\!\mathrm{LN}\big(\mathbf{g}_i+\mathrm{GELU}(W\!\sum_j A_{ij}\mathbf{g}_j)\big)$, $\times2$};

% dashed group encloses only the two gating blocks
\begin{scope}[on background layer]
  \node[grp, draw=boxcol, line width=1.6pt, fit=(gate)(kl)] (gbox) {};
\end{scope}

% fused row
\node[blbl, anchor=east] (frow) at (25.65,-2.32) {$A_{i,:}=\sum_k \pi_{ik}A^{k}_{i,:}$};
\foreach \c [count=\q from 0] in {white,teal!40,purple!55,orange!45,teal!25}
  \node[cell, fill=\c, draw=black!75, line width=0.9pt] at (24.9+0.55*\q,-2.8) {};
\foreach \t [count=\q from 0] in {i,j,k,l,m}
  \node[font=\normalsize, text=black!60] at (24.9+0.55*\q,-3.3) {$\t$};
\node[blbl, anchor=north, text width=60mm, align=center] (rs) at (26.0,-3.65)
  {row-stochastic (Lem.~1)};

% arrows: gate -> KL -> message passing ; fused row -> message passing
\draw[f] (gate.south) -- (kl.north);
\draw[f, draw=purple!60, dashed] (kl.south) -- (mp.north);
\draw[F] (26.0,-2.5) -- (mp.south);
\draw[f] (1.55,-4.85) -- (26.0,-4.85) -- (26.0,-4.35);
\end{tikzpicture}%
}
\caption{Component~1 of Fig.~\ref{fig:arch}: multi-graph fusion, shown for one focal region $i$ (shaded). Top: the same five regions carry four adjacencies, a fixed distance kernel, a learned static graph, a state-dependent attention graph, and the genomic transmission-linkage graph, whose edge widths are proportional to the number of near-identical genome pairs. Bottom: row $i$ of each adjacency, shaded by weight, is scaled by that region's gate coefficient $\pi_{ik}$ and summed into a single row-stochastic row. The genomic row is masked where a region has no sequenced cases ($m$), and the gate renormalises over the remaining channels.}
\label{fig:fusion}
\end{figure*}
